\chapter{Requisiti: elicitazione e analisi}

Nel Capitolo 3 abbiamo incontrato la prima fase del modello a cascata, la \textbf{Requirements Engineering}, e abbiamo visto
la regola d'oro: in questa fase si decide \textit{cosa} il sistema deve fare, non \textit{come} lo farà.
Ora entriamo nel dettaglio: come si scoprono i requisiti, come si scrivono e come si evita che siano ambigui o inutilizzabili.

\begin{center}
\begin{tcolorbox}[enhanced, breakable=false, colback=primary!5, colframe=primary!70!black, boxrule=1pt, arc=3pt,
  width=0.92\textwidth, drop fuzzy shadow, fontupper=\itshape]
  \iquotel\ La parte più difficile nella costruzione di un sistema software è decidere con precisione cosa costruire.
  Nessun'altra parte del lavoro concettuale è così difficile, nessun'altra danneggia tanto il sistema finale se fatta male,
  nessun'altra è più difficile da correggere in seguito.\\[4pt]
  \normalfont\hfill\textbf{--- Fred P. Brooks} (parafrasi)
\end{tcolorbox}
\end{center}

\dfn{Requirements Engineering (RE)}{
  Sotto-area dell'Ingegneria del Software che si occupa del processo di \textbf{definizione dei requisiti} per il software
  da realizzare. Il suo scopo è fornire agli ingegneri \textbf{metodi, tecniche e strumenti} per capire e documentare
  che cosa un sistema software deve fare.
}

% ============================================================
\section{Cos'è un requisito}

\dfn{Requisito}{
  Un \textbf{requisito} è la descrizione di ciò che il sistema deve fare: i \textbf{servizi} che deve fornire ai suoi utenti
  e i \textbf{vincoli operativi} entro cui deve funzionare.
}

Qualche esempio, riferito a una piattaforma per la gestione dei corsi universitari:

\begin{table}[H]
\centering
\renewcommand{\arraystretch}{1.35}
\begin{tabularx}{\textwidth}{L >{\centering\arraybackslash}p{3.6cm}}
\toprule
Requisito & Tipo \\
\midrule
Il sistema deve permettere agli studenti di cercare i corsi per parole chiave. & servizio \\
Il sistema deve permettere agli studenti di iscriversi a un corso. & servizio \\
Le interfacce utente devono essere realizzate come app native Android e iOS. & vincolo \\
Il sistema deve gestire almeno 10.000 studenti collegati contemporaneamente. & vincolo \\
Il sistema deve autenticare gli utenti anche tramite il sensore di impronte digitali del dispositivo. & servizio + vincolo \\
\bottomrule
\end{tabularx}
\end{table}

\subsection{Due livelli: requisiti utente e requisiti di sistema}

Nell'industria il termine «requisito» viene usato in modo non uniforme: a volte indica una descrizione generale e astratta,
altre volte una definizione molto dettagliata. Per questo si distinguono due \textbf{livelli}.

\dfn{Requisiti utente (User Requirements)}{
  Descrizione \textbf{astratta e di alto livello} di un servizio che il sistema deve fornire o di un vincolo che deve rispettare.
  Il punto di vista è quello dell'\textbf{utente finale}, quindi sono scritti in modo comprensibile a chi non è tecnico.
}

\dfn{Requisiti di sistema (System Requirements)}{
  Definizione \textbf{dettagliata e formale} di una funzione del sistema. Il punto di vista è quello del \textbf{sistema da costruire},
  quindi sono precisi e non lasciano spazio a interpretazioni.
}

\begin{table}[H]
\centering
\renewcommand{\arraystretch}{1.4}
\begin{tabularx}{\textwidth}{L L}
\toprule
\textcolor{ph1}{\textbf{Requisito utente}} & \textcolor{ph4}{\textbf{Requisiti di sistema corrispondenti}} \\
\midrule
\textbf{1.} Il sistema deve generare un report mensile delle iscrizioni ai corsi. &
\textbf{1.1} I report delle iscrizioni devono essere generati l'ultimo giorno lavorativo di ogni mese.\newline
\textbf{1.2} Deve essere generato un report per ogni corso, che mostri il numero totale di iscrizioni del mese corrente
e l'età media degli studenti iscritti.\newline
\textbf{1.3} L'accesso ai report deve essere riservato ai responsabili dei corsi.\newline
\textbf{1.4} \ldots \\
\bottomrule
\end{tabularx}
\caption{Un requisito utente si espande in molti requisiti di sistema.}
\end{table}

Perché mantenere due livelli invece di uno solo? Perché servono a due scopi diversi, entrambi di natura \textbf{contrattuale}.

\begin{figure}[H]
\centering
\begin{tikzpicture}
  \node[card=ph1, text width=4.4cm, minimum height=2.5cm, align=left] (ur) {\textbf{Requisiti utente}\\[3pt]\scriptsize
     Il cliente definisce, a grandi linee, ciò che gli serve.\\[3pt]
     Base per il \textbf{bando di gara}: aziende diverse possono proporre soluzioni diverse.};
  \node[card=swgray, text width=3.2cm, minimum height=2.5cm, align=center, right=0.8cm of ur] (bid) {\textbf{Gara}\\[3pt]\scriptsize Le aziende presentano le loro offerte.\\[3pt] Una vince il contratto.};
  \node[card=ph4, text width=4.4cm, minimum height=2.5cm, align=left, right=0.8cm of bid] (sr) {\textbf{Requisiti di sistema}\\[3pt]\scriptsize
     L'azienda scelta li formula in dettaglio, così il cliente capisce esattamente cosa riceverà e può validarli.};
  \node[card=swred, text width=3.2cm, minimum height=2.5cm, align=center, right=0.8cm of sr] (ct) {\textbf{Contratto}\\[3pt]\scriptsize Una volta accettati, i requisiti di sistema diventano \textbf{vincolanti}.};
  \draw[flow] (ur) -- (bid);
  \draw[flow] (bid) -- (sr);
  \draw[flow] (sr) -- (ct);
\end{tikzpicture}
\caption{I due livelli di requisiti hanno una doppia funzione contrattuale.}
\end{figure}

\warning{Perché i requisiti valgono tanto}{
  Un errore nei requisiti è l'errore più costoso che si possa fare. Come abbiamo visto nel Capitolo 3 con la curva di Boehm,
  un difetto individuato in fase di requisiti costa \textbf{1}; lo stesso difetto scoperto quando il sistema è già in esercizio
  può costare \textbf{500 volte tanto}. In più, se i requisiti di sistema sono già nel contratto, ogni modifica diventa
  anche una trattativa legale ed economica.
}

% ============================================================
\section{Requisiti funzionali e non funzionali}

La classificazione più usata divide i requisiti in due famiglie.

\begin{center}
\begin{tikzpicture}
  \node[card=ph1, text width=7.9cm, align=left, minimum height=3.6cm] (f) {%
    \textbf{\large Requisiti funzionali}\\[4pt]\small
    Descrivono i \textbf{servizi} che il sistema deve fornire, come deve \textbf{reagire a particolari input}
    e come deve comportarsi in determinate situazioni.\\[4pt]
    \scriptsize\textit{Es.}: «Un utente deve potersi iscrivere a un corso.»\\
    \scriptsize\textit{Es.}: «Un utente deve poter annullare la propria iscrizione a un corso.»};
  \node[card=ph3, text width=7.9cm, align=left, minimum height=3.6cm, right=0.4cm of f] (nf) {%
    \textbf{\large Requisiti non funzionali}\\[4pt]\small
    Sono \textbf{vincoli} sui servizi o sulle funzioni offerte: vincoli di tempo, vincoli sul processo di sviluppo,
    vincoli imposti da standard. Spesso riguardano il sistema \textbf{nel suo insieme}, non la singola funzione.\\[4pt]
    \scriptsize\textit{Es.}: «Il sistema deve essere conforme al GDPR.»\\
    \scriptsize\textit{Es.}: «Il sistema deve gestire fino a 10.000 utenti contemporanei.»};
\end{tikzpicture}
\end{center}

\begin{esercizio}{Funzionale o non funzionale?}
Classifica i seguenti requisiti.
\begin{enumerate}
  \item Un utente deve potersi iscrivere a un corso.
  \item Il sistema deve essere conforme al GDPR.
  \item Il sistema deve gestire fino a 10.000 utenti contemporanei.
  \item Un utente deve poter cercare i corsi che contengono certe parole chiave.
  \item Un utente deve poter annullare la propria iscrizione a un corso.
  \item Il sistema deve funzionare su una macchina con 8 GB di memoria.
\end{enumerate}

\textbf{Soluzione.} (1) F \quad (2) NF \quad (3) NF \quad (4) F \quad (5) F \quad (6) NF.
\end{esercizio}

\warning{La distinzione non è sempre netta}{
  Consideriamo il requisito: «solo gli utenti autorizzati devono poter accedere al sistema». Sembra non funzionale
  (è un vincolo di sicurezza sull'intero sistema), ma se lo sviluppiamo in dettaglio genera requisiti chiaramente \textbf{funzionali}:
  «l'utente deve poter effettuare il login e autenticarsi», «il sistema deve mostrare un messaggio di errore dopo tre tentativi falliti»\ldots

  Morale: i requisiti \textbf{non sono indipendenti}: un requisito spesso ne genera o ne vincola altri.
}

\subsection{Come si scrivono le due famiglie}

\begin{itemize}
  \item I requisiti \textbf{funzionali utente} si scrivono in \textbf{linguaggio naturale}, in modo che siano comprensibili anche
        a persone non tecniche (utenti, dirigenti).
  \item I requisiti \textbf{funzionali di sistema} devono essere \textbf{dettagliati}: descrivono con precisione input, output
        ed eccezioni, così che gli sviluppatori sappiano esattamente cosa implementare.
\end{itemize}

I requisiti non funzionali, invece, non riguardano un servizio specifico: vincolano \textbf{come} il sistema deve essere realizzato
(ad esempio «deve essere sviluppato in Java») o specificano proprietà globali (tempi di risposta, uso della memoria, \ldots).

\error{I requisiti non funzionali sono spesso i più critici}{
  Se un requisito funzionale è implementato male, gli utenti spesso trovano un modo per aggirare il problema.
  Se invece \textbf{non} è soddisfatto un requisito non funzionale, il sistema può diventare \textbf{inutilizzabile}:
  un sistema non conforme al GDPR non può proprio essere messo in esercizio in Europa, per quanto bene funzioni.
}

\subsection{Tipi di requisiti non funzionali}

I requisiti non funzionali possono nascere da tre sorgenti diverse: dal \textbf{prodotto}, dall'\textbf{organizzazione}
(del cliente o di chi sviluppa) o da \textbf{fonti esterne}.

\begin{figure}[H]
\centering
\begin{tikzpicture}[font=\scriptsize,
  figlio/.style 2 args={subchip=#1, minimum width=3.9cm, minimum height=0.6cm, anchor=west, text width=3.5cm, align=center},
  nipote/.style 2 args={subchip=#1, fill=#1!28, minimum width=2.6cm, minimum height=0.55cm, anchor=west}]

  \node[chip=ph3, minimum width=4.6cm, minimum height=0.8cm, font=\small\bfseries] (root) at (0,0) {Requisiti non funzionali};

  % --- le tre famiglie
  \node[chip=q3, minimum width=4.4cm, minimum height=0.7cm, font=\scriptsize\bfseries] (prod) at (-5.8,-1.5) {Requisiti di prodotto};
  \node[chip=q7, minimum width=4.4cm, minimum height=0.7cm, font=\scriptsize\bfseries] (org)  at (0,-1.5)    {Requisiti organizzativi};
  \node[chip=q5, minimum width=4.4cm, minimum height=0.7cm, font=\scriptsize\bfseries] (ext)  at (5.8,-1.5)  {Requisiti esterni};
  \foreach \n in {prod,org,ext}{ \draw[line width=0.9pt, swgray] (root.south) -- ++(0,-0.45) -| (\n.north); }

  % --- colonna PRODOTTO (binario a x=-8.25, figli con anchor west a x=-7.95)
  \node[figlio={q3}{}] (fu) at (-7.95,-2.35) {Usabilità};
  \node[figlio={q3}{}] (fe) at (-7.95,-3.07) {Efficienza};
  \node[figlio={q3}{}] (fa) at (-7.95,-5.09) {Affidabilità};
  \node[figlio={q3}{}] (fs) at (-7.95,-5.81) {Sicurezza (security)};
  \draw[line width=0.8pt, swgray] (-8.25,-1.85) -- (-8.25,-5.81);
  \foreach \n in {fu,fe,fa,fs}{ \draw[line width=0.8pt, swgray] (-8.25,-1.85) |- (\n.west); }
  % sotto-figli di Efficienza (binario indentato a x=-7.3)
  \node[nipote={q3}{}] (np) at (-7.0,-3.75) {Prestazioni};
  \node[nipote={q3}{}] (ns) at (-7.0,-4.37) {Spazio};
  \draw[line width=0.7pt, swgray!80] (-7.3,-3.32) -- (-7.3,-4.37);
  \foreach \n in {np,ns}{ \draw[line width=0.7pt, swgray!80] (-7.3,-3.37) |- (\n.west); }

  % --- colonna ORGANIZZATIVI
  \node[figlio={q7}{}] (oa) at (-2.15,-2.35) {Ambientali};
  \node[figlio={q7}{}] (oo) at (-2.15,-3.07) {Operativi};
  \node[figlio={q7}{}] (os) at (-2.15,-3.79) {Di sviluppo};
  \draw[line width=0.8pt, swgray] (-2.45,-1.85) -- (-2.45,-3.79);
  \foreach \n in {oa,oo,os}{ \draw[line width=0.8pt, swgray] (-2.45,-1.85) |- (\n.west); }

  % --- colonna ESTERNI
  \node[figlio={q5}{}] (er) at (3.65,-2.35) {Regolamentari};
  \node[figlio={q5}{}] (ee) at (3.65,-3.07) {Etici};
  \node[figlio={q5}{}] (el) at (3.65,-3.79) {Legislativi};
  \draw[line width=0.8pt, swgray] (3.35,-1.85) -- (3.35,-3.79);
  \foreach \n in {er,ee,el}{ \draw[line width=0.8pt, swgray] (3.35,-1.85) |- (\n.west); }
  % sotto-figli di Legislativi
  \node[nipote={q5}{}] (ec) at (4.6,-4.47) {Contabili};
  \node[nipote={q5}{}] (es) at (4.6,-5.09) {Di sicurezza};
  \draw[line width=0.7pt, swgray!80] (4.3,-4.04) -- (4.3,-5.09);
  \foreach \n in {ec,es}{ \draw[line width=0.7pt, swgray!80] (4.3,-4.09) |- (\n.west); }
\end{tikzpicture}
\caption{Tassonomia dei requisiti non funzionali in base alla loro origine.}
\end{figure}

\begin{table}[H]
\centering
\renewcommand{\arraystretch}{1.4}
\begin{tabularx}{\textwidth}{>{\bfseries}P{3.4cm} L L}
\toprule
Origine & Cosa riguardano & Esempi \\
\midrule
\textcolor{q3}{Di prodotto} & Caratteristiche richieste al prodotto stesso & Velocità del sistema, memoria occupata, usabilità, tasso di guasti accettabile \\
\textcolor{q7}{Organizzativi} & Politiche e procedure dell'organizzazione del cliente o dello sviluppatore & Processi operativi con cui il sistema sarà usato, linguaggio di programmazione obbligatorio, ambiente operativo da rispettare \\
\textcolor{q5}{Esterni} & Fattori esterni al sistema e al processo di sviluppo & Approvazione da parte di un ente regolatore, obblighi di legge, requisiti etici per l'accettabilità da parte di utenti e pubblico \\
\bottomrule
\end{tabularx}
\end{table}

\subsection{Requisiti di dominio}

\dfn{Requisiti di dominio}{
  Requisiti che derivano dallo \textbf{specifico settore o contesto industriale} in cui il software viene usato.
  Possono essere sia funzionali sia non funzionali.
}

\begin{esempio}{Un requisito di dominio per un dispositivo medico}
«La sicurezza del sistema deve essere garantita secondo lo standard IEC 60601-1: \textit{Medical Electrical Equipment --- Part 1:
General Requirements for Basic Safety and Essential Performance}.»

Questo requisito vincola sia la \textbf{progettazione} del sistema sia il \textbf{processo di sviluppo}: un software perfettamente
funzionante ma non sviluppato secondo lo standard \textbf{non potrebbe essere usato} in pratica.
\end{esempio}

\info{Misurare i requisiti non funzionali}{
  Per essere utile, un requisito non funzionale deve essere associato a una \textbf{metrica}:

  \begin{center}
  \renewcommand{\arraystretch}{1.25}
  \begin{tabular}{>{\bfseries}l l}
  \toprule
  Proprietà & Metriche tipiche \\
  \midrule
  Prestazioni & Operazioni elaborate al secondo; tempo di risposta; tempo di aggiornamento dello schermo \\
  Dimensione & Megabyte occupati \\
  Facilità d'uso & Ore di formazione necessarie; tasso di errore degli utenti; richieste di assistenza \\
  Affidabilità & Tempo medio tra due guasti; percentuale di disponibilità \\
  Robustezza & Tempo di ripristino dopo un guasto; probabilità di perdita di dati \\
  \bottomrule
  \end{tabular}
  \end{center}
}

% ============================================================
\section{Le proprietà di un buon requisito}

\begin{table}[H]
\centering
\renewcommand{\arraystretch}{1.45}
\begin{tabularx}{\textwidth}{>{\bfseries\leavevmode\color{ph1}}P{3.2cm} L}
\toprule
Proprietà & Significato \\
\midrule
Chiaro & Facile da capire, soprattutto per i requisiti utente. \\
Non ambiguo & Una sola interpretazione possibile: l'ambiguità porta a \textbf{contenziosi} con il cliente. \\
Completo & Non deve mancare nulla di ciò che serve. \\
Coerente & I requisiti non devono contraddirsi tra loro. \\
Testabile & Dato un requisito, deve essere possibile stabilire \textbf{senza dubbi} se il sistema lo soddisfa. \\
Tracciabile & Deve essere collegabile «in avanti e all'indietro» alle altre parti della documentazione: progetto, codice, test. \\
\bottomrule
\end{tabularx}
\end{table}

\subsection{I demoni dell'ambiguità}

Il linguaggio naturale è comodo ma insidioso. Consideriamo un cartello stradale che tutti capiamo al volo:

\begin{figure}[H]
\centering
\begin{minipage}[c]{0.24\textwidth}
\centering
\begin{tikzpicture}[scale=0.95]
  % cartello "school" (pentagono giallo)
  \fill[yellow-gradient-6!85!black] (-1.0,2.1) -- (1.0,2.1) -- (1.0,3.3) -- (0,3.95) -- (-1.0,3.3) -- cycle;
  \draw[draw=black, line width=1pt] (-1.0,2.1) -- (1.0,2.1) -- (1.0,3.3) -- (0,3.95) -- (-1.0,3.3) -- cycle;
  \node[font=\bfseries\scriptsize, text=black] at (0,2.85) {SCHOOL};
  % cartello del limite di velocita'
  \draw[draw=black, fill=white, line width=1pt] (-1.05,-0.75) rectangle (1.05,1.85);
  \node[font=\tiny\bfseries, text=black, align=center] at (0,1.55) {SPEED\\LIMIT};
  \node[font=\Huge\bfseries, text=black] at (0,0.75) {20};
  \node[font=\tiny\bfseries, text=black, align=center] at (0,-0.3) {WHEN\\CHILDREN\\ARE PRESENT};
  \draw[swgray, line width=2pt] (0,-0.75) -- (0,-1.35);
\end{tikzpicture}
\end{minipage}\hfill
\begin{minipage}[c]{0.72\textwidth}
Per un essere umano l'intento è chiaro: proteggere i bambini in una zona scolastica. Ma proviamo a spiegarlo a un computer:
\begin{itemize}
  \item perché il cartello dice «children» (plurale) e non «child»? La regola vale solo se c'è \textbf{più di un} bambino?
  \item chi è considerato un «bambino»? Un ragazzo di 17 anni lo è?
  \item cosa significa essere «presenti»? Dentro la scuola? Sul marciapiede? Visibili dalla strada?
  \item la regola vale solo quando la scuola è aperta?
\end{itemize}
\end{minipage}
\caption{Un cartello chiarissimo per una persona, pieno di ambiguità per una macchina (H. Robinson, \textit{The demons of ambiguity}).}
\end{figure}

\begin{esempio}{Le domande nascoste in un requisito}
Requisito: \textit{«L'applicazione deve permettere agli utenti di avviare e arrestare un servizio su una macchina remota.»}

Sembra chiaro, ma:
\begin{itemize}
  \item cosa succede se si avvia un servizio \textbf{già in esecuzione}? E se si arresta un servizio \textbf{non in esecuzione}?
  \item se la macchina remota si \textbf{riavvia}, il servizio deve ripartire automaticamente?
  \item se cade la \textbf{connessione}: l'applicazione mostra un messaggio d'errore? Riprova? Dopo quanto tempo?
        Quanti tentativi al massimo?
\end{itemize}
Ognuna di queste domande, se non trova risposta nei requisiti, diventerà una discussione (o un contenzioso) con il cliente.
\end{esempio}

\subsection{Testabilità}

Deve essere possibile stabilire \textbf{senza ambiguità} se il sistema soddisfa un requisito: altrimenti gli sviluppatori
sosterranno che il requisito è soddisfatto e il cliente sosterrà il contrario. In pratica:

\begin{itemize}
  \item i requisiti \textbf{funzionali di sistema} devono essere il più dettagliati possibile e non lasciare spazio a interpretazioni;
  \item i requisiti \textbf{non funzionali} devono includere, quando possibile, \textbf{indicatori quantitativi}.
\end{itemize}

\begin{table}[H]
\centering
\renewcommand{\arraystretch}{1.45}
\begin{tabularx}{\textwidth}{L L}
\toprule
\textcolor{swred}{\textbf{\icross\ Poco testabile}} & \textcolor{swgreen}{\textbf{\itick\ Testabile}} \\
\midrule
Il sistema deve essere affidabile. & Il sistema deve avere una disponibilità mensile del 99,99\%. \\
Il sistema deve essere facile da usare. & Dopo al più due ore di formazione, gli utenti devono saper usare tutte le funzioni del sistema; dopo la formazione, il numero medio di errori non deve superare i due per ora di utilizzo. \\
Il sistema deve essere veloce e reattivo. & Il tempo medio di risposta del sistema non deve superare i 100 ms. \\
\bottomrule
\end{tabularx}
\caption{Trasformare requisiti vaghi in requisiti verificabili.}
\end{table}

% ============================================================
\section{Il processo di Requirements Engineering}

Il processo di RE si articola in tre attività principali.

\begin{figure}[H]
\centering
\begin{tikzpicture}
  \node[phase=ph1, minimum width=4.2cm, minimum height=1.2cm, font=\small\bfseries, drop shadow] (el) at (0,0) {Elicitazione\\e analisi};
  \node[phase=ph3, minimum width=4.2cm, minimum height=1.2cm, font=\small\bfseries, drop shadow] (sp) at (5.9,0) {Specifica};
  \node[phase=ph5, minimum width=4.2cm, minimum height=1.2cm, font=\small\bfseries, drop shadow] (va) at (11.8,0) {Validazione};
  \draw[flow] (el) -- (sp);
  \draw[flow] (sp) -- (va);
  \draw[feedback] (va.north) to[out=120, in=60] node[above, font=\scriptsize, text=swred]{i problemi trovati fanno tornare indietro} (el.north);
  \node[card=ph1, text width=4.2cm, font=\scriptsize, align=left, below=0.5cm of el] (d1) {\textbf{Scopo}: scoprire i requisiti interagendo con gli stakeholder.\\[3pt]\textit{Output}: descrizione del sistema.};
  \node[card=ph3, text width=4.2cm, font=\scriptsize, align=left, below=0.5cm of sp] (d2) {\textbf{Scopo}: convertire i requisiti in una forma standardizzata.\\[3pt]\textit{Output}: requisiti utente e di sistema.};
  \node[card=ph5, text width=4.2cm, font=\scriptsize, align=left, below=0.5cm of va] (d3) {\textbf{Scopo}: controllare che i requisiti definiscano davvero il sistema voluto dal cliente.\\[3pt]\textit{Output}: documento dei requisiti.};
\end{tikzpicture}
\caption{Le tre attività del processo di Requirements Engineering.}
\end{figure}

\subsection{Un processo iterativo, non lineare}

Nella pratica queste attività non si susseguono una sola volta: si \textbf{intrecciano} in un processo iterativo che lavora
a \textbf{livelli di granularità} diversi.

\begin{enumerate}
  \item Nella prima iterazione ci si concentra sui \textbf{requisiti di business}: perché il progetto è necessario?
        Chi ne beneficia? Quando e dove verrà usato?
  \item Poi si definiscono i \textbf{requisiti utente}.
  \item Infine si definiscono i \textbf{requisiti di sistema}.
\end{enumerate}

Ogni iterazione comprende elicitazione e analisi, specifica e validazione.

\begin{figure}[H]
\centering
\begin{tikzpicture}[font=\scriptsize]
  \foreach \c/\tit/\x in {0/{Elicitazione / Analisi}/0, 1/{Specifica}/5.3, 2/{Validazione}/10.6}{
    \node[font=\small\bfseries, text=swgray] at (\x,1.1) {\tit};
  }
  \foreach \r/\liv/\col/\val/\y in {0/{Business}/q1/{Studio di fattibilità}/0,
                                    1/{Utente}/q4/{Prototipazione}/-1.8,
                                    2/{Sistema}/q2/{Revisione}/-3.6}{
    \node[phase=\col, minimum width=4.4cm, minimum height=1.1cm, font=\scriptsize\bfseries] (a\r) at (0,\y) {Elicitazione e analisi\\dei requisiti \liv};
    \node[phase=\col, minimum width=4.4cm, minimum height=1.1cm, font=\scriptsize\bfseries] (b\r) at (5.3,\y) {Specifica dei requisiti\\\liv};
    \node[phase=\col, fill=\col!70, minimum width=4.4cm, minimum height=1.1cm, font=\scriptsize\bfseries] (c\r) at (10.6,\y) {\val};
    \node[card=\col, minimum width=2.4cm, minimum height=1.1cm, font=\scriptsize\bfseries, anchor=west] (d\r) at (13.3,\y) {Requisiti\\\liv};
    \draw[flow] (a\r) -- (b\r);
    \draw[flow] (b\r) -- (c\r);
    \draw[flow] (c\r) -- (d\r);
  }
  \foreach \r [evaluate=\r as \s using int(\r+1)] in {0,1}{
    \draw[-{Stealth[length=2.5mm]}, line width=1pt, swgray] (d\r.south) -- ++(0,-0.35) -| ([xshift=-0.55cm]a\s.north west) -- ([xshift=-0.02cm]a\s.west);
  }
\end{tikzpicture}
\caption{Il processo iterativo: ogni livello di requisiti attraversa le tre attività e alimenta il livello successivo.}
\end{figure}

% ============================================================
\section{Elicitazione e analisi dei requisiti}

\dfn{Elicitazione dei requisiti}{
  Attività il cui obiettivo è \textbf{scoprire che cosa i clienti vogliono e di cosa hanno bisogno}, interagendo con gli
  stakeholder. È considerata la parte \textbf{più critica} dell'intera ingegneria dei requisiti.
}

Detta così sembra banale; in realtà è la parte più difficile del mestiere, per un motivo semplice: le persone non sanno
sempre cosa vogliono, e quasi mai lo sanno spiegare in termini utilizzabili da chi deve costruire il software.

\subsection{Chi sono gli stakeholder}

\dfn{Stakeholder}{
  Le \textbf{persone o i gruppi} che sono in qualche modo interessati o influenzati dal progetto software.
}

\begin{figure}[H]
\centering
\begin{tikzpicture}
  \node[circle, fill=ph1!12, draw=ph1, line width=1.2pt, minimum size=2.6cm, align=center, font=\small\bfseries, text=ph1!80!black] (sys) at (0,0) {Sistema\\da costruire};
  \foreach \ang/\tit/\desc/\col in {
      120/{Utenti finali}/{gruppo \textbf{eterogeneo}: chi inserisce gli ordini (vendite), chi li evade (logistica), chi assiste i clienti}/q4,
      60/{Chi usa gli output}/{non usa il sistema, ma lavora sui dati che produce (report, statistiche)}/q3,
      -60/{Dirigenza del cliente}/{decide, paga e valuta i benefici per l'organizzazione}/q1,
      -120/{Personale tecnico}/{dovrà installare, configurare e far funzionare il sistema ogni giorno}/q7}{
    \node[card=\col, text width=4.6cm, font=\scriptsize, align=left] (n\ang) at (\ang:5.2) {\textbf{\iuser\ \tit}\\[2pt]\desc};
    \draw[line width=1pt, \col] (sys) -- (n\ang);
  }
\end{tikzpicture}
\caption{Gli stakeholder tipici di un sistema di gestione degli ordini.}
\end{figure}

\subsection{Perché è difficile}

\begin{figure}[H]
\centering
\begin{tikzpicture}
  \foreach \i/\tit/\txt/\col in {
    0/{Non sanno cosa vogliono}/{Gli stakeholder hanno idee solo generali. Non sanno cosa è tecnicamente fattibile e possono fare richieste irrealistiche.}/swred,
    1/{Parlano il loro gergo}/{Sono esperti del \textbf{loro} dominio ed esprimono i requisiti con i loro termini, dando per scontate conoscenze implicite.}/sworange,
    2/{Dicono cose diverse}/{Stakeholder diversi esprimono requisiti diversi e in modi diversi: bisogna trovare i punti in comune e risolvere i conflitti.}/q1}{
    \node[card=\col, text width=5.2cm, minimum height=3.0cm, align=left, font=\scriptsize, anchor=north west] at (\i*5.85,0) {\textbf{\normalsize\tit}\\[4pt]\txt};
  }
  \foreach \i/\tit/\txt/\col in {
    0/{Ci sono fattori politici}/{Un dirigente può chiedere una funzionalità perché gli fa aumentare la propria influenza nell'organizzazione, non perché serva.}/q7,
    1/{Le priorità sono in conflitto}/{Chi sente che le proprie esigenze non sono state considerate può arrivare a \textbf{ostacolare} il processo.}/q5,
    2/{I requisiti cambiano}/{Durante l'elicitazione emergono nuovi requisiti, magari da stakeholder che all'inizio non erano stati considerati.}/q2}{
    \node[card=\col, text width=5.2cm, minimum height=3.0cm, align=left, font=\scriptsize, anchor=north west] at (\i*5.85,-3.4) {\textbf{\normalsize\tit}\\[4pt]\txt};
  }
\end{tikzpicture}
\caption{Le sei principali difficoltà dell'elicitazione dei requisiti.}
\end{figure}

\subsection{Il processo di elicitazione e analisi}

L'attività si svolge come un \textbf{ciclo} che si ripete finché i requisiti non sono sufficientemente stabili.

\begin{figure}[H]
\centering
\begin{tikzpicture}
  \node[phase=ph1, minimum width=3.9cm, minimum height=1.1cm, font=\scriptsize\bfseries] (d1) at (0,2.3) {1. Scoperta e\\comprensione};
  \node[phase=ph2, minimum width=3.9cm, minimum height=1.1cm, font=\scriptsize\bfseries] (d2) at (4.3,0) {2. Classificazione\\e organizzazione};
  \node[phase=ph4, minimum width=3.9cm, minimum height=1.1cm, font=\scriptsize\bfseries] (d3) at (0,-2.3) {3. Prioritizzazione\\e negoziazione};
  \node[phase=ph6, minimum width=3.9cm, minimum height=1.1cm, font=\scriptsize\bfseries] (d4) at (-4.3,0) {4. Documentazione};
  \node[font=\scriptsize\itshape, text=swgray, align=center] at (0,0) {il ciclo si ripete\\finché i requisiti\\non sono stabili};
  \draw[flow, line width=1.4pt] (d1) to[out=0, in=90] (d2);
  \draw[flow, line width=1.4pt] (d2) to[out=270, in=0] (d3);
  \draw[flow, line width=1.4pt] (d3) to[out=180, in=270] (d4);
  \draw[flow, line width=1.4pt] (d4) to[out=90, in=180] (d1);
\end{tikzpicture}
\caption{Il ciclo di elicitazione e analisi dei requisiti.}
\end{figure}

\begin{table}[H]
\centering
\renewcommand{\arraystretch}{1.35}
\begin{tabularx}{\textwidth}{>{\bfseries}P{5.0cm} L}
\toprule
\textcolor{ph1}{1. Scoperta e comprensione} & Si interagisce con gli stakeholder per \textbf{scoprire} i requisiti. \\
\textcolor{ph2}{2. Classificazione e organizzazione} & I requisiti correlati vengono \textbf{raggruppati} e organizzati in categorie coerenti. \\
\textcolor{ph4}{3. Prioritizzazione e negoziazione} & Si risolvono i \textbf{conflitti} che nascono dalle esigenze contrastanti dei diversi stakeholder. \\
\textcolor{ph6}{4. Documentazione} & I requisiti trovati vengono \textbf{documentati}, come base per l'iterazione successiva. \\
\bottomrule
\end{tabularx}
\end{table}

% ============================================================
\section{Le tecniche di elicitazione}

Non esiste una tecnica unica: si combinano approcci diversi.

\begin{center}
\begin{tikzpicture}
  \foreach \i/\tit/\ico/\col in {0/{Interviste}/\iuser/ph1, 1/{Osservazione\\(etnografia)}/\ieye/ph2, 2/{Personas}/\ibadge/ph3,
                                 3/{Storie e scenari}/\ibook/ph4, 4/{Mock-up e\\prototipi}/\iwindow/ph5}{
    \node[card=\col, text width=2.9cm, minimum height=1.7cm, align=center, font=\small\bfseries] at (\i*3.4,0) {{\Large\ico}\\[3pt]\tit};
  }
\end{tikzpicture}
\end{center}

% ------------------------------------------------------------
\subsection{Interviste con gli stakeholder}

Le interviste, formali o informali, sono lo strumento fondamentale della RE: il team chiede agli stakeholder come lavorano,
quale sistema usano oggi e cosa si aspettano dal sistema da sviluppare. I requisiti si ricavano dalle risposte.

\begin{table}[H]
\centering
\renewcommand{\arraystretch}{1.4}
\begin{tabularx}{\textwidth}{>{\bfseries}P{3.4cm} L}
\toprule
Intervista chiusa & Lo stakeholder risponde a una \textbf{lista di domande predefinita}. \\
Intervista aperta & Non c'è un'agenda prestabilita: si lascia parlare l'intervistato. \\
\midrule
\textcolor{swgreen}{\textbf{In pratica}} & Si usa una \textbf{combinazione} delle due. Le interviste completamente aperte funzionano raramente: conviene partire da un insieme di domande preparate, che di solito portano ad altri temi da discutere in modo meno strutturato. \\
\bottomrule
\end{tabularx}
\end{table}

\begin{center}
\begin{tikzpicture}
  \node[card=swgreen, text width=7.7cm, minimum height=2.6cm, align=left, font=\small] (pro) {%
    \textbf{\itick\ Il lato buono}\\[3pt]\scriptsize
    Alle persone in genere piace parlare del proprio lavoro e partecipano alle interviste con un atteggiamento positivo.};
  \node[card=swred, text width=7.7cm, minimum height=2.6cm, align=left, font=\small, right=0.4cm of pro] (con) {%
    \textbf{\icross\ Le difficoltà}\\[3pt]\scriptsize
    \textbullet\ usano \textbf{gergo} tecnico e danno per scontati molti requisiti;\\
    \textbullet\ conoscono bene il \textbf{proprio} lavoro, ma hanno una visione parziale o distorta di quello dei colleghi di altri reparti;\\
    \textbullet\ possono essere \textbf{reticenti} su vincoli organizzativi, per via di equilibri di potere interni.};
\end{tikzpicture}
\end{center}

\warning{Rispettare il tempo degli stakeholder}{
  Gli stakeholder hanno il loro lavoro da fare: le interviste vanno preparate bene, per non sprecare il loro tempo
  (e il nostro) con riunioni inutili.
}

\begin{esempio}{Il gergo: una storia vera}
Durante un tirocinio presso la direzione finanziaria di un grande Comune, un tributo minore veniva gestito quasi interamente
con fogli di calcolo. Alla domanda «mi racconta come funziona la gestione di questo tributo?», la risposta è stata:

\begin{center}
\begin{tcolorbox}[enhanced, breakable=false, colback=swteal!8, colframe=swteal, boxrule=0.8pt, arc=3pt, width=0.95\linewidth, fontupper=\itshape\small]
  «Certo! È molto semplice! Un operatore sul campo mi invia la \textbf{lettura} per un dato anno per un certo contribuente.
  Quindi inserisco la lettura in una maschera che calcola l'\textbf{imponibile}. Successivamente faccio la stampa delle
  \textbf{ingiunzioni di pagamento}, che vengono poi emesse. Una volta notificata l'ingiunzione, il contribuente ha 60 giorni
  per pagare. Se non paga entro il termine, devo procedere all'\textbf{iscrizione al ruolo coattivo}: inserisco i dati in una
  maschera che genera il \textbf{Tracciato 290} da trasmettere all'ente di riscossione.»
\end{tcolorbox}
\end{center}

«È molto semplice» --- ma per capire davvero quella frase servono settimane di studio del dominio. Questo è esattamente il
problema del gergo: per l'intervistato è ovvio, per l'ingegnere del software è una lingua straniera. Il lavoro consiste nel
farsi spiegare ogni termine, finché la descrizione non diventa un insieme di requisiti precisi.
\end{esempio}

% ------------------------------------------------------------
\subsection{Etnografia}

\dfn{Etnografia}{
  Tecnica \textbf{osservativa} in cui l'ingegnere dei requisiti si \textbf{immerge nell'ambiente di lavoro} degli utenti finali:
  osserva il lavoro quotidiano e prende nota dei compiti realmente svolti e delle persone coinvolte.
}

Il software non vive isolato: è usato in un ambiente \textbf{sociale e organizzativo} che può generare o vincolare requisiti.
E, molto spesso, il modo in cui si lavora davvero è \textbf{diverso} dal processo formale che si dovrebbe seguire.

L'etnografia è particolarmente efficace per scoprire:
\begin{itemize}
  \item requisiti che derivano dal modo in cui le persone \textbf{lavorano realmente}, invece che dal modo in cui le procedure
        aziendali dicono che dovrebbero lavorare;
  \item requisiti che derivano dalla \textbf{cooperazione} e dalla consapevolezza delle attività altrui.
\end{itemize}

\begin{esempio}{Un requisito che nessuno avrebbe raccontato}
Osservando l'ufficio, si nota che prima di inviare un questionario di soddisfazione al cliente, un addetto del servizio clienti
telefona ogni volta a un collega della logistica per sapere se l'ordine è stato davvero consegnato.

Nessuno lo avrebbe mai detto in un'intervista («si fa così e basta»), ma da lì nasce un requisito preciso:
\textit{«Gli addetti al servizio clienti devono poter verificare se un cliente ha già ricevuto il proprio ordine.»}
\end{esempio}

% ------------------------------------------------------------
\subsection{Personas}

\dfn{Persona}{
  Un \textbf{archetipo ipotetico} di utente reale. Le \textit{personas} non sono persone vere, ma le rappresentano:
  sono immaginarie, eppure definite con \textbf{rigore e precisione}, e non sono inventate di sana pianta, ma \textbf{scoperte}
  come sottoprodotto del processo di elicitazione.
}

Conoscere l'utente è fondamentale sia per scoprire i requisiti sia (lo vedremo più avanti) per progettare interfacce efficaci.
A volte gli utenti li conosciamo bene (ad esempio se sviluppiamo uno strumento per altri ingegneri del software), ma in generale
il software è fatto \textbf{per altre persone}, con esigenze diverse dalle nostre: la nostra esperienza personale non basta
a immaginare come lavorano e cosa vogliono.

Le personas servono anche quando \textbf{non abbiamo stakeholder da intervistare}: ad esempio quando sviluppiamo un prodotto
generico per il grande pubblico, o quando parte degli utenti finali è il pubblico generico.

Non esiste un formato standard; di solito una persona contiene:

\begin{center}
\begin{tikzpicture}
  \foreach \i/\t/\d in {0/{Personalizzazione}/{nome, età, breve biografia},
                        1/{Lavoro}/{che lavoro fa?},
                        2/{Istruzione}/{studi ed esperienza},
                        3/{Obiettivi}/{cosa si aspetta dal software?},
                        4/{Frustrazioni}/{quali problemi il software può risolverle?}}{
    \node[card=ph3, text width=2.8cm, minimum height=1.6cm, align=center, font=\scriptsize] at (\i*3.25,0) {\textbf{\t}\\[3pt]\d};
  }
\end{tikzpicture}
\end{center}

\begin{esempio}{Una persona per un gestionale di studio dentistico}
Supponiamo di raccogliere i requisiti per un software di gestione degli appuntamenti di uno studio dentistico.
Intervistando il personale abbiamo già ricavato alcuni requisiti:
\begin{itemize}\itemsep1pt
  \item i pazienti devono potersi registrare e autenticare nel sistema;
  \item i pazienti devono poter prenotare un appuntamento in uno slot libero;
  \item i medici devono poter vedere gli appuntamenti della settimana corrente;
  \item i receptionist devono poter cancellare gli appuntamenti;
  \item i pazienti devono ricevere un promemoria via e-mail il giorno prima dell'appuntamento;
  \item i pazienti devono ricevere un'e-mail se il loro appuntamento viene annullato.
\end{itemize}

Costruiamo ora una persona:

\begin{center}
\begin{tcolorbox}[enhanced, breakable=false, colback=white, colframe=ph3, boxrule=1pt, arc=3pt, drop fuzzy shadow,
   width=\linewidth, left=6pt, right=6pt]
\begin{minipage}[t]{0.30\linewidth}
  \centering
  \begin{tikzpicture}
    \fill[ph3!15, rounded corners=4pt] (-1.3,-1.5) rectangle (1.3,1.5);
    \fill[ph3!55] (0,0.55) circle (0.45);
    \fill[ph3!55] (-0.75,-1.2) .. controls (-0.75,0.05) and (0.75,0.05) .. (0.75,-1.2) -- cycle;
  \end{tikzpicture}\\[4pt]
  {\small\textbf{Concetta (Tina) Aiello}}\\[2pt]
  {\scriptsize Età: 68 --- vedova\\ Pensionata (ex maestra elementare)\\ \textit{assertiva, impaziente}}
\end{minipage}\hfill
\begin{minipage}[t]{0.66\linewidth}
  \small\textbf{Bio.} Tina è un'insegnante di scuola primaria in pensione, con oltre 40 anni di esperienza. È dolce e paziente
  con i bambini, ma diventa impaziente quando si trova fuori dalla sua zona di comfort. Ama la sua routine quotidiana e non
  le piace cambiarla. La tecnologia non le piace: l'anno scorso i figli le hanno regalato il primo smartphone, che usa solo per
  videochiamare le nipoti. È diffidente verso internet, che percepisce come un posto pericoloso.\\[3pt]
  \textbf{Interessi.} Lavoro a maglia; comunità della parrocchia; tempo con le nipoti.\\[3pt]
  \textbf{Obiettivi.} Prenotare i controlli periodici della protesi dentale; ricevere un promemoria per non dimenticare
  l'appuntamento; spostare un appuntamento se coincide con un altro impegno.
\end{minipage}
\end{tcolorbox}
\end{center}

\textbf{Cosa abbiamo scoperto.} Tina non prenoterà mai da sola tramite l'app: telefonerà allo studio. E un promemoria via e-mail
non lo leggerà. Nascono quindi \textbf{nuovi requisiti} che nessuno aveva pensato di dirci:
\begin{itemize}\itemsep1pt
  \item i receptionist devono poter creare appuntamenti per i pazienti che telefonano allo studio;
  \item il giorno prima dell'appuntamento, i receptionist devono ricevere una notifica che ricordi loro di \textbf{telefonare}
        ai pazienti che hanno prenotato per telefono.
\end{itemize}
\end{esempio}

% ------------------------------------------------------------
\subsection{Storie e scenari}

\dfn{Storia (story)}{
  Testo \textbf{narrativo} che presenta una descrizione ad alto livello dell'uso del sistema. È efficace per comunicare
  gli obiettivi «di insieme» in modo comprensibile a chiunque.
}

\begin{center}
\begin{tcolorbox}[enhanced, breakable=false, colback=ph4!5, colframe=ph4, boxrule=0.9pt, arc=3pt,
  title={\textbf{Storia: prenotare un appuntamento con l'app}}, coltitle=white, colbacktitle=ph4, fontupper=\small]
  Pamela ha bisogno di prenotare una visita odontoiatrica perché il dente del giudizio ha iniziato a farle male.
  Accede al sistema e cerca gli slot disponibili per oggi o per domani. Il sistema mostra l'elenco degli slot liberi.
  Pamela ne sceglie uno e conferma la prenotazione. Riceve un'e-mail di conferma con i dettagli dell'appuntamento.
\end{tcolorbox}
\end{center}

\dfn{Scenario}{
  Descrizione (tipicamente \textbf{strutturata}) di un esempio di interazione tra utente e sistema. Raffina una storia
  aggiungendo informazioni precise: stato iniziale, flusso normale degli eventi, cosa può andare storto e come si gestisce,
  altre attività in corso nello stesso momento, stato del sistema alla fine.
}

\begin{table}[H]
\centering
\renewcommand{\arraystretch}{1.4}
\begin{tabularx}{\textwidth}{>{\bfseries\leavevmode\color{ph4}}P{3.4cm} L}
\toprule
\multicolumn{2}{l}{\textcolor{ph4}{\textbf{Scenario: prenotazione di un appuntamento}}} \\
\midrule
Assunzioni iniziali & L'utente ha un account sul sistema di gestione degli appuntamenti. \\
Flusso normale & L'utente effettua il login e cerca gli slot disponibili per la data odierna o per il giorno successivo.
Il sistema mostra l'elenco degli slot liberi nelle date scelte. L'utente seleziona uno degli slot proposti.
Il sistema invia un'e-mail di conferma all'utente, al receptionist e al medico. \\
Cosa può andare storto & \textbullet\ Non c'è nessuno slot disponibile nelle date considerate: il sistema mostra un messaggio
di errore e chiede di scegliere date diverse.\newline
\textbullet\ Lo slot scelto non è più disponibile al momento della conferma, perché un altro utente lo ha prenotato:
il sistema mostra un messaggio di errore e torna alla schermata degli slot disponibili. \\
Stato finale del sistema & Il sistema ha memorizzato il nuovo appuntamento. \\
\bottomrule
\end{tabularx}
\caption{Uno scenario strutturato che raffina la storia precedente.}
\end{table}

% ------------------------------------------------------------
\subsection{Mock-up a bassa fedeltà}

\dfn{Mock-up a bassa fedeltà (wireframe)}{
  \textbf{Schizzi semplificati} delle interfacce utente del sistema. Sono un modo rapido ed economico per visualizzare
  la struttura e i flussi principali: l'attenzione è sulla \textbf{funzionalità e sul flusso}, non sull'estetica.
}

\begin{figure}[H]
\centering
\begin{minipage}[c]{0.46\textwidth}
\centering
\begin{tikzpicture}[line width=0.9pt]
  % schermata 1
  \draw[swgray, rounded corners=3pt] (0,0) rectangle (4.2,5.4);
  \draw[swgray] (0,4.7) -- (4.2,4.7);
  \node[font=\scriptsize\bfseries, text=swgray] at (2.1,5.05) {Prenota una visita};
  \node[font=\tiny, text=swgray, anchor=west] at (0.25,4.25) {Data};
  \draw[swgray!70] (0.25,3.5) rectangle (3.95,4.0);
  \node[font=\tiny, text=swgray!80, anchor=west] at (0.35,3.75) {gg / mm / aaaa};
  \node[font=\tiny, text=swgray, anchor=west] at (0.25,3.15) {Slot disponibili};
  \foreach \k in {0,1,2}{
    \draw[swgray!70] (0.25,2.25-\k*0.65) rectangle (3.95,2.8-\k*0.65);
    \node[font=\tiny, text=swgray!90, anchor=west] at (0.4,2.52-\k*0.65) {0\the\numexpr9+\k\relax:00 --- Dott. Bianchi};
    \draw[swgray!70] (3.35,2.35-\k*0.65) rectangle (3.85,2.7-\k*0.65);
  }
  \draw[fill=swgray!25, swgray, rounded corners=2pt] (1.1,0.35) rectangle (3.1,0.9);
  \node[font=\tiny\bfseries, text=swgray] at (2.1,0.62) {CONFERMA};
\end{tikzpicture}
\end{minipage}\hfill
\begin{minipage}[c]{0.50\textwidth}
\textbf{Perché servono}
\begin{itemize}\itemsep2pt
  \item \textbf{Facilitano la comprensione}: aiutano gli stakeholder a cogliere le funzionalità di base e sono un punto di
        partenza per la discussione.
  \item \textbf{Fanno scoprire nuovi requisiti}: è molto più facile ragionare su un'interfaccia concreta che su frasi astratte.
  \item \textbf{Coinvolgono e permettono di iterare}: si ottiene feedback rapido senza investire in grafica.
  \item \textbf{Sono una base per lo sviluppo}: da qui si passa ai prototipi ad alta fedeltà e poi all'implementazione.
\end{itemize}
\end{minipage}
\caption{Un wireframe non deve essere bello: deve far capire struttura e flusso.}
\end{figure}

\info{Strumenti}{
  \textbf{Carta e matita}: rapidi e accessibili per i primi schizzi, ma diventano disordinati per sistemi complessi.\\
  \textbf{Strumenti digitali}: software commerciali come \textit{Balsamiq} o \textit{Figma} permettono di creare e condividere
  wireframe elettronicamente; esistono anche alternative open source come \textit{mydraft.cc}, \textit{draw.io} o \textit{Pencil Project}.
}

% ============================================================
\section{La specifica dei requisiti}

\dfn{Specifica dei requisiti}{
  Il processo di \textbf{mettere per iscritto} i requisiti utente e di sistema in un \textbf{documento dei requisiti}.
  I requisiti utente devono essere comprensibili a utenti e clienti senza background tecnico; i requisiti di sistema sono
  più dettagliati e possono contenere informazioni tecniche.
}

Poiché il documento dei requisiti può diventare \textbf{parte del contratto} di sviluppo, è importante che sia il più
completo e dettagliato possibile.

\warning{Requisiti e progetto sono inseparabili}{
  In linea di principio i requisiti dicono \textit{cosa} il sistema deve fare e il progetto dice \textit{come} lo fa.
  In pratica i due aspetti si intrecciano:
  \begin{itemize}
    \item i requisiti possono essere \textbf{organizzati} seguendo un'architettura di sistema di alto livello;
    \item il sistema può dover \textbf{interoperare} con altri sistemi, e questo genera requisiti di progetto;
    \item l'uso di una specifica architettura per soddisfare requisiti non funzionali può essere esso stesso un requisito di dominio.
  \end{itemize}
}

\subsection{Quattro approcci alla specifica}

\begin{figure}[H]
\centering
\begin{tikzpicture}
  \draw[-{Stealth[length=3mm]}, line width=1.2pt, swgray] (-0.4,-0.9) -- (15.6,-0.9) node[right, font=\scriptsize, align=left]{};
  \node[font=\scriptsize, text=swgray] at (7.6,-1.35) {formalità e precisione crescenti $\longrightarrow$};
  \foreach \i/\tit/\txt/\col in {
    0/{Linguaggio naturale}/{Frasi numerate in linguaggio naturale: una frase, un requisito.}/q4,
    1/{Linguaggio naturale strutturato}/{Si usa un \textbf{modulo standard} o un template per ogni requisito.}/q1,
    2/{Notazioni semi-formali}/{Modelli UML (es.\ diagrammi dei casi d'uso) e modelli di dominio, con annotazioni testuali.}/q2,
    3/{Specifica formale}/{Notazioni basate su concetti matematici: automi a stati finiti, logiche temporali.}/q7}{
    \node[card=\col, text width=3.3cm, minimum height=2.7cm, align=left, font=\scriptsize, anchor=south west] at (\i*4.0,-0.7) {\textbf{\small\tit}\\[4pt]\txt};
  }
\end{tikzpicture}
\caption{Dal testo libero alla matematica: quattro modi di specificare i requisiti.}
\end{figure}

Approcci diversi si adattano a domini diversi:

\begin{table}[H]
\centering
\renewcommand{\arraystretch}{1.4}
\begin{tabularx}{\textwidth}{>{\bfseries}P{5.2cm} L}
\toprule
Contesto & Approccio tipico \\
\midrule
Sistemi critici per la sicurezza & Specifiche \textbf{formali} e linguaggio naturale strutturato \\
Una semplice app per le liste di cose da fare & Linguaggio naturale non strutturato \\
Un sistema informativo di media o grande dimensione & Notazioni \textbf{semi-formali} e modelli: un buon compromesso \\
\bottomrule
\end{tabularx}
\end{table}

\subsection{Specifica in linguaggio naturale}

I requisiti sono scritti come frasi in linguaggio naturale, eventualmente accompagnate da diagrammi e tabelle.
Si usa perché è \textbf{espressivo, intuitivo e universale}: i requisiti risultano comprensibili a utenti e clienti,
e permette di esprimere sia requisiti funzionali sia non funzionali.

\begin{center}
\begin{tikzpicture}
  \node[card=swgreen, text width=7.7cm, minimum height=3.5cm, align=left, font=\scriptsize] (g) {%
    \textbf{\normalsize\itick\ Linee guida}\\[4pt]
    \textbullet\ definire un \textbf{formato standard} e usarlo per tutti i requisiti, ad esempio:\\
    \hspace*{0.3cm}\textit{<parte del sistema> shall <requisito> (<motivazione>)};\\
    \textbullet\ usare il linguaggio in modo \textbf{coerente}: \textit{shall} per i requisiti obbligatori,
    \textit{should} per quelli desiderabili;\\
    \textbullet\ \textbf{evidenziare} graficamente le parti chiave del requisito;\\
    \textbullet\ evitare il \textbf{gergo informatico};\\
    \textbullet\ includere sempre la \textbf{motivazione} (\textit{rationale}) del requisito.};
  \node[card=swred, text width=7.7cm, minimum height=3.5cm, align=left, font=\scriptsize, right=0.4cm of g] (b) {%
    \textbf{\normalsize\icross\ Problemi tipici}\\[4pt]
    \textbullet\ \textbf{Mancanza di chiarezza}: è difficile essere precisi senza rendere il documento pesante da leggere;\\
    \textbullet\ \textbf{confusione tra requisiti}: requisiti funzionali e non funzionali tendono a mescolarsi;\\
    \textbullet\ \textbf{amalgama di requisiti}: più requisiti diversi finiscono espressi insieme in un'unica frase.};
\end{tikzpicture}
\end{center}

\begin{esempio}{Un requisito con la sua motivazione}
\textit{«Il report mensile complessivo delle iscrizioni \textbf{shall} includere statistiche sul genere degli studenti iscritti
(necessarie per compilare il report sulla parità di genere).»}

La motivazione tra parentesi è preziosa: se un domani quell'obbligo dovesse cadere, sapremo che quel requisito può essere rimosso.
\end{esempio}

\subsection{Specifica strutturata}

\dfn{Specifica strutturata}{
  Approccio in cui la \textbf{libertà} di chi scrive i requisiti viene limitata e si impone un \textbf{modo standardizzato}
  di scriverli (un modulo con campi prestabiliti).
}

Funziona bene per alcuni tipi di requisiti (ad esempio nei sistemi di controllo embedded), mentre risulta troppo rigida per i
requisiti dei sistemi gestionali.

\begin{table}[H]
\centering
\renewcommand{\arraystretch}{1.35}
\begin{tabularx}{\textwidth}{>{\bfseries\leavevmode\color{ph2}}P{2.8cm} L}
\toprule
\multicolumn{2}{l}{\textcolor{ph2}{\textbf{Microinfusore di insulina / Software di controllo / SRS / 3.3.2}}} \\
\midrule
Funzione & Calcolo della dose di insulina con livello di zucchero nella norma. \\
Descrizione & Calcola la dose di insulina da somministrare quando il livello di zucchero misurato è compreso tra 3 e 7 unità. \\
Input & Lettura corrente dello zucchero ($r_2$); le due letture precedenti ($r_0$ e $r_1$). \\
Sorgente & Lettura corrente dal sensore; le altre letture dalla memoria. \\
Output & \texttt{CompDose}: la dose di insulina da somministrare. \\
Azione & \texttt{CompDose} è zero se il livello è stabile o in calo, oppure se sta crescendo ma con velocità di crescita
decrescente. Se il livello cresce e la velocità di crescita aumenta, \texttt{CompDose} si calcola dividendo per 4 la differenza
tra il livello corrente e quello precedente e arrotondando. Se il risultato arrotondato è zero, \texttt{CompDose} è posta
alla dose minima somministrabile. \\
Richiede & Le due letture precedenti, per poter calcolare la velocità di variazione. \\
Pre-condizione & Il serbatoio contiene almeno la dose singola massima ammessa. \\
Post-condizione & In memoria, $r_0$ è sostituito da $r_1$ e poi $r_1$ è sostituito da $r_2$. \\
Effetti collaterali & Nessuno. \\
\bottomrule
\end{tabularx}
\caption{Esempio di specifica strutturata (da I. Sommerville).}
\end{table}

\subsection{Come si scrive una frase di requisito}

Lo standard ISO/IEC/IEEE 29148 detta regole molto precise sul linguaggio da usare:

\begin{table}[H]
\centering
\renewcommand{\arraystretch}{1.4}
\begin{tabularx}{\textwidth}{>{\bfseries\leavevmode\color{ph1}}P{2.6cm} L}
\toprule
shall & Introduce un \textbf{requisito obbligatorio e vincolante}. È la forma da usare per i requisiti veri e propri. \\
will & Introduce \textbf{dichiarazioni di fatto}, di intenzione o di contesto: non è vincolante. \\
should & Introduce \textbf{preferenze o obiettivi} desiderabili ma non obbligatori: non sono requisiti. \\
may & Introduce \textbf{suggerimenti o concessioni}: non è vincolante. \\
must & \textbf{Da evitare}: rischia di essere interpretato erroneamente come requisito. \\
\midrule
\multicolumn{2}{l}{\textbf{Altre regole di scrittura}} \\
\midrule
Forma positiva & Usare affermazioni positive ed evitare requisiti negativi come «\textit{shall not}». \\
Forma attiva & Usare la forma attiva, evitando il passivo («\textit{it is required that\ldots}»). \\
Verbi precisi & Evitare espressioni come «\textit{shall be able to}»: appesantiscono senza aggiungere precisione. \\
\bottomrule
\end{tabularx}
\end{table}

\error{Termini vaghi da evitare}{
  Le parole vaghe generano requisiti impossibili da verificare o interpretabili in più modi. In particolare:
  \begin{itemize}\itemsep1pt
    \item \textbf{superlativi}: «il migliore», «il massimo»;
    \item \textbf{linguaggio soggettivo}: «intuitivo», «facile da usare», «conveniente»;
    \item \textbf{pronomi vaghi}: «esso», «questo», «quello»;
    \item \textbf{avverbi e aggettivi ambigui}: «quasi sempre», «significativo», «minimo»; e congiunzioni ambigue come «e/o»;
    \item \textbf{termini aperti e non verificabili}: «fornire supporto», «incluso ma non limitato a», «come minimo»;
    \item \textbf{frasi comparative}: «migliore di», «di qualità superiore»;
    \item \textbf{scappatoie}: «se possibile», «se appropriato», «ove applicabile»;
    \item \textbf{termini assoluti}: «tutti», «sempre», «mai», «ogni» (difficilissimi da verificare);
    \item \textbf{riferimenti incompleti}: citare un documento senza data, versione o sezione applicabile.
  \end{itemize}
}

\subsection{Un modello per i requisiti funzionali}

Un modo pratico per scrivere requisiti funzionali chiari è seguire un \textbf{pattern} fisso:

\begin{figure}[H]
\centering
\begin{tikzpicture}
  \foreach \i/\t/\col in {0/{[Condizione]}/q3, 1/{[Soggetto]}/q1, 2/{[Azione]}/ph1, 3/{[Oggetto]}/q7, 4/{[Vincolo dell'azione]}/q5}{
    \node[chip=\col, minimum width=3.0cm, minimum height=0.85cm, font=\small\bfseries] at (\i*3.15,0) {\t};
  }
\end{tikzpicture}
\end{figure}

\begin{table}[H]
\centering
\renewcommand{\arraystretch}{1.45}
\begin{tabularx}{\textwidth}{>{\bfseries}P{3.2cm} L}
\toprule
Forma completa & \textit{Quando viene ricevuto il segnale x} [Condizione], \textit{il sistema} [Soggetto] \textit{deve impostare}
[Azione] \textit{il bit di ricezione del segnale x} [Oggetto] \textit{entro 2 secondi} [Vincolo]. \\
Variante & \textit{Con stato del mare pari a 1} [Condizione], \textit{il sistema radar} [Soggetto] \textit{deve rilevare}
[Azione] \textit{i bersagli} [Oggetto] \textit{fino a 100 miglia nautiche di distanza} [Vincolo]. \\
Senza condizione & \textit{Il sistema di fatturazione} [Soggetto] \textit{deve mostrare le fatture dei clienti in sospeso}
[Azione] \textit{in ordine crescente di data di scadenza} [Vincolo]. \\
\bottomrule
\end{tabularx}
\caption{Lo stesso schema applicato a tre requisiti diversi.}
\end{table}

% ============================================================
\section{I documenti dei requisiti}

Esistono standard che definiscono la \textbf{struttura} del documento dei requisiti.

\subsection{IEEE 830-1998}

Definisce una struttura generica da adattare a ogni specifico sistema:

\begin{table}[H]
\centering
\renewcommand{\arraystretch}{1.3}
\begin{tabularx}{\textwidth}{>{\bfseries\leavevmode\color{ph2}}P{4.2cm} L}
\toprule
1. Introduzione & Scopo del documento; ambito del prodotto; definizioni, acronimi e abbreviazioni; riferimenti; panoramica dell'organizzazione del documento. \\
2. Descrizione generale & Prospettiva del prodotto (in quale contesto sarà usato? interagirà con sistemi esistenti?); funzioni del prodotto; caratteristiche degli utenti; vincoli generali (leggi, regolamenti); assunzioni e dipendenze (hardware o software esistenti). \\
3. Requisiti specifici & Requisiti funzionali (per ciascuno: introduzione, input, elaborazione, output); requisiti di interfaccia esterna (utente, hardware, software, comunicazioni); requisiti di prestazione; vincoli di progetto (conformità a standard, limiti hardware); attributi del sistema software (sicurezza, manutenibilità); altri requisiti (database, operatività, adattamento al sito). \\
4. Appendici & Materiale di supporto. \\
5. Indice analitico & \\
\bottomrule
\end{tabularx}
\caption{Struttura di un documento dei requisiti secondo lo standard IEEE 830-1998.}
\end{table}

\subsection{ISO/IEC/IEEE 29148:2018}

Nel 2011 è stato pubblicato un nuovo standard, aggiornato nel 2018: \textit{Systems and software engineering --- Life cycle
processes requirements engineering}. Non definisce un unico documento, ma \textbf{quattro documenti} corrispondenti ai diversi
livelli visti all'inizio del capitolo.

\begin{figure}[H]
\centering
\begin{tikzpicture}
  \foreach \i/\sig/\nome/\desc/\col in {
     0/{BRS}/{Business Requirements Specification}/{Perché il sistema viene sviluppato o modificato: processi, politiche e regole dell'organizzazione, requisiti di alto livello dal punto di vista degli stakeholder. È la base per la partecipazione attiva degli stakeholder al processo.}/q1,
     1/{StRS}/{Stakeholder Requirements Specification}/{Nel contesto descritto dal BRS, descrive come l'organizzazione userà il sistema per contribuire al business, esprimendo le esigenze di utenti, operatori e manutentori derivate dal contesto d'uso.}/q4,
     2/{SyRS}/{System Requirements Specification}/{Requisiti tecnici del sistema e usabilità dell'interazione uomo-sistema: obiettivi generali, ambiente di destinazione, vincoli, assunzioni e requisiti non funzionali. Descrive cosa il sistema deve fare in termini di interazioni con l'esterno.}/q2,
     3/{SRS}/{Software Requirements Specification}/{Specifica di un particolare prodotto software: le funzioni che deve svolgere in uno specifico ambiente. Può essere scritta da rappresentanti del fornitore, del committente o di entrambi.}/q7}{
    \node[chip=\col, minimum width=2.0cm, minimum height=0.9cm, font=\small\bfseries, anchor=north west] (s\i) at (0.6,-\i*2.3) {\sig};
    \node[card=\col, text width=12.9cm, minimum height=1.7cm, align=left, font=\scriptsize, anchor=north west] (c\i) at (3.0,-\i*2.3+0.1) {\textbf{\small\nome}\\[3pt]\desc};
    \draw[line width=0.8pt, swgray!70] (s\i.east) -- (c\i.west);
  }
  % freccia di livello, a sinistra dei riquadri
  \draw[-{Stealth[length=3.5mm]}, line width=1.6pt, swgray] (0.0,-0.15) -- (0.0,-7.25);
  \node[font=\scriptsize\bfseries, text=swgray, rotate=90, anchor=south] at (-0.15,-3.7) {dal business al software};
\end{tikzpicture}
\caption{I quattro documenti previsti dallo standard ISO/IEC/IEEE 29148:2018.}
\end{figure}

\info{Cosa deve specificare la parte «Functions»}{
  Lo standard entra molto nel dettaglio. Per le funzioni del software richiede, ad esempio, di definire:
  i controlli di validità sugli input; l'esatta sequenza delle operazioni; le risposte a situazioni anomale
  (overflow, problemi di comunicazione, guasti hardware, gestione e recupero degli errori); l'effetto dei parametri;
  la relazione tra output e input (sequenze di input/output e formule di conversione).
}

% ============================================================
\section{Esercizi}

\begin{esercizio}{Rendere testabili i requisiti}
Riscrivi i seguenti requisiti in forma verificabile.
\begin{enumerate}
  \item «Il sistema deve avere un'interfaccia intuitiva.»
  \item «Il sistema deve gestire un gran numero di utenti.»
  \item «I dati devono essere salvati in modo sicuro.»
\end{enumerate}

\textbf{Soluzione (una delle possibili).}
(1) «Un utente senza formazione precedente deve completare la prenotazione di un appuntamento in meno di 3 minuti, senza consultare
il manuale, nell'80\% dei casi.»
(2) «Il sistema deve servire almeno 5.000 utenti collegati contemporaneamente con un tempo medio di risposta inferiore a 500 ms.»
(3) «Tutti i dati personali devono essere memorizzati cifrati con algoritmo AES-256 e accessibili solo agli utenti con ruolo
\textit{amministratore}.»
\end{esercizio}

\begin{esercizio}{Trova l'ambiguità}
Individua i problemi dei seguenti requisiti.
\begin{enumerate}
  \item «Il sistema deve inviare una notifica all'utente appena possibile dopo la registrazione e, se appropriato, anche al suo referente.»
  \item «Il modulo di reportistica deve essere migliore di quello attuale e supportare tutti i formati.»
\end{enumerate}

\textbf{Soluzione.}
(1) Contiene la scappatoia «se appropriato» (chi decide?), il termine aperto «appena possibile» (entro quanto?) e
\textbf{amalgama} due requisiti distinti (notifica all'utente e notifica al referente) in un'unica frase.
(2) Usa una \textbf{frase comparativa} non verificabile («migliore di») e un \textbf{termine assoluto} («tutti i formati»)
che è impossibile da dimostrare: andrebbero elencati i formati richiesti.
\end{esercizio}

\begin{esercizio}{Quale tecnica di elicitazione?}
Indica la tecnica più adatta a ciascuna situazione.
\begin{enumerate}
  \item Si sospetta che gli operatori del magazzino seguano una procedura diversa da quella ufficiale.
  \item Si sviluppa un'app per il grande pubblico e non esiste un cliente specifico da intervistare.
  \item Si vuole capire come il responsabile della segreteria gestisce oggi le iscrizioni.
  \item Si vuole far discutere gli stakeholder su come sarà organizzata la schermata principale.
\end{enumerate}

\textbf{Soluzione.} (1) Etnografia (osservazione sul campo). (2) Personas. (3) Intervista (semi-strutturata).
(4) Mock-up a bassa fedeltà.
\end{esercizio}

\gbox{In sintesi}{azure-gradient-6!80!black}{
\begin{itemize}
  \item Un \textbf{requisito} descrive un servizio che il sistema deve offrire o un vincolo che deve rispettare.
        Si distinguono \textbf{requisiti utente} (astratti, per il cliente) e \textbf{requisiti di sistema} (dettagliati, vincolanti in contratto).
  \item I requisiti si dividono in \textbf{funzionali} (cosa fa il sistema) e \textbf{non funzionali} (vincoli sul sistema nel suo insieme:
        di prodotto, organizzativi, esterni). I requisiti \textbf{di dominio} derivano dal settore d'uso.
  \item Un buon requisito è \textbf{chiaro, non ambiguo, completo, coerente, testabile e tracciabile};
        i requisiti non funzionali vanno espressi con \textbf{indicatori quantitativi}.
  \item Il processo di RE comprende \textbf{elicitazione e analisi}, \textbf{specifica} e \textbf{validazione}, ripetute in modo
        \textbf{iterativo} sui livelli business, utente e sistema.
  \item Le tecniche di elicitazione principali sono \textbf{interviste}, \textbf{etnografia}, \textbf{personas},
        \textbf{storie e scenari}, \textbf{mock-up a bassa fedeltà}.
  \item La specifica può essere in linguaggio naturale, in linguaggio naturale \textbf{strutturato}, con notazioni
        \textbf{semi-formali} (UML) o \textbf{formali}. Lo standard ISO/IEC/IEEE 29148 regola persino l'uso di
        \textit{shall}, \textit{should}, \textit{will} e \textit{may}.
\end{itemize}
}
